Extracting brain tumor boundaries from \gls{mri} accurately is an important task in medical image analysis, since it affects how surgery is planned and how nearby tissue is evaluated. Zotin et al. proposed an automated pipeline for this problem \cite{zotin2018edge}, organized into three phases. First, an image enhancement stage uses a median filter for noise suppression and \gls{bcet} to improve contrast in the region of interest. Second, a segmentation stage applies \gls{fcm} clustering to divide the image into homogeneous classes by pixel intensity, followed by the Canny operator to extract thin edges of the tumor. Finally, a data-analysis stage evaluates the result with metrics such as \gls{fom}, Sensitivity, and Accuracy, comparing the predicted edges to reference maps drawn by human experts \cite{pratt1979quantitative}. The original study reports good results on a small, hand-picked set of images, but unsupervised intensity-based methods do not necessarily generalize to the larger anatomical variability found in clinical datasets. This motivates further pre-processing, such as morphological skull-stripping \cite{hooda2014brain} and deterministic centroid initialization for \gls{fcm} \cite{benson2016brain, szilagyi2003mri}. The rest of this paper is organized as follows. Section \ref{sec:baseline} reproduces the baseline pipeline as closely as possible to the original paper and validates it on a much larger dataset, showing where it struggles. Section \ref{sec:optimized} proposes an optimized pipeline that addresses these issues through dynamic slice selection, a switch to \gls{flair} sequences, and a more tolerant evaluation metric. Section \ref{sec:comparison} compares the baseline and optimized results and discusses the role of \gls{csf} interference. Section \ref{sec:conclusions} summarizes the findings.